She did not sleep, which surprised no one who knew her, least of all herself.

She sat instead in the command tent with a single brazier burning down to embers and did the arithmetic, the way Grey Worm had known she would, the way she had done arithmetic her whole life over every wound the world had ever dealt her — turning the thing over and over, examining it from every angle, refusing herself the mercy of not looking directly at it.

\textit{Aegon Targaryen.} Her brother Rhaegar’s son. The boy the Usurper’s rebellion had supposedly killed in his cradle, dashed against a wall by the Mountain’s own hands — except that boy had died, truly, at King’s Landing, and this was a different son, a hidden son, a son born in a tower in the south to a Stark girl everyone had called a kidnap victim and a Targaryen prince everyone had called her abductor. A son raised his whole life as a bastard, beneath a false name, by a man who had loved him well enough to spend his own honor hiding him.

She found, examining it, that she believed it. That was the first surprise. She had expected to want to disbelieve it, to reach for the comfort of calling it a lie, Sam Tarly’s grief-poisoned invention aimed at wounding her in the one place a queen could always be wounded. But it fit too well. It explained too much — the way Jon fought, the way the dragons had taken to him, Rhaegal accepting his weight without complaint when no living creature but a Targaryen had ever ridden a dragon and lived. She had told herself, these past weeks, that it was simply the bond of a good man and a good beast. She understood now it had been blood calling to blood, and she had been too content in her own certainty to hear it.

\sceneBreak

So. The man she loved had a better claim to her throne than she did.

She made herself say it plainly, in the privacy of her own skull, because she had learned long ago that a truth left unspoken grew larger and stranger in the dark than it ever was in the light. Jon Snow — Aegon Targaryen — trueborn, male, elder in the line by the laws every lord in Westeros recognized. If he stood in a hall and claimed it, and if the claim could be proven, then everything she had crossed the world to take was his by right, and hers only by the accident of his having been hidden.

And he had known. For weeks. He had knelt beside her in the dark and shared her bed and her counsel and her war, and known the whole time that he stood ahead of her, and said nothing.

That was the wound. Not the claim itself — she found, testing it, that the claim frightened her less than she’d expected. It was the silence. The weeks of it. The knowledge that a man she had trusted more completely than she trusted anyone living had carried the single most dangerous truth in Westeros in his chest and looked at her across a hundred fires and chosen, every time, not to give it to her.

She turned that over a long while, and found, slowly, that it admitted of two readings, and that she could not make herself choose between them.

The first reading was the kind one. \textit{He said nothing because the claim means nothing to him.} Jon Snow had given up a crown once already, had bent his knee to her freely when he might have stayed King in the North, had shown her again and again that he wanted the throne about as much as he wanted a wasting sickness. A man like that hid the truth not to hoard a weapon but to keep from ever having to hold it — because telling her would have forced a choice he did not want to make, between a claim he never asked for and an alliance he believed in. In this reading his silence was not betrayal at all. It was a kind of loyalty, clumsy and frightened, but real.

The second reading was the cold one. \textit{He said nothing because a hidden claim is worth more than a spent one.} A truth kept secret was a truth a man could produce at the moment of his choosing — when the Northern lords who already resented her had grown resentful enough, when the war had weakened her enough, when the moment came to decide who Westeros would actually kneel to once the dead were beaten. In this reading Jon’s silence was the most dangerous kind of patience, and his talk of not wanting the throne only the sort of thing a wise claimant said right up until the day the claiming became convenient.

She did not believe the second reading. That was the thing she kept returning to, in the long cold hours. She \textit{did not believe it.} She knew Jon Snow, knew the particular incorruptible dullness of his honor, knew he would sooner fall on his own sword than scheme for a crown behind a lover’s back.

And yet she could not make herself fully disbelieve it either, and that — the not-being-able, the small cold seed of doubt that would not quite die however many times she reasoned it dead — was what told her, near dawn, that something between them had broken that would not mend back to what it was, whatever she decided to do about it.

\sceneBreak

She decided to do nothing.

She arrived at it slowly, the way she arrived at all her hardest decisions now, and understood, arriving, exactly what it cost and exactly why she was choosing it regardless.

To confront him would force the thing into the open, and the open was the one place this truth could not survive being. A hall full of starving Northern lords, already half-mutinous over grain and a foreign queen, told that the man they’d crowned themselves was the trueborn heir to the very throne she’d crossed the world to claim — that was not a conversation. That was a war inside a war, and the living could not afford a second war with the dead still marching at their backs. Whatever Jon’s silence had been, kind or cold, he had at least been right about that much: this truth told now would not be a gift. It would be a blade, and it would cut the alliance in half at the exact moment the alliance was the only thing keeping any of them alive.

So she would say nothing. She would let Jon go on believing his secret was still his own. She would keep the alliance whole because the alliance had to be whole, and she would carry the knowledge the way he had carried it, in silence, across a hundred more fires.

But she would watch him now. She could not help that, and did not try to. She would sit in council and love him and plan a war beside him, and some cold, sleepless part of her would be watching the whole time, weighing his words, listening for the moment — if it ever came, and she prayed it never would — when the kind reading and the cold reading finally forced her to choose between them.

She had told Tyrion, only that morning, that she no longer entirely trusted herself to know where the line was anymore. She understood now, watching the brazier gutter toward dawn, that she had just found a new one, and stepped across it, and that no one in the world knew she had done it, and that this was perhaps the loneliest thing that had ever happened to her — to carry a wound this large entirely alone, and choose, for the good of everyone else, to let no one see that she was bleeding at all.